\section{Conclusiones}
Fue verificada la reutilización de una habilitación periódica para controlar dos comportamientos secuenciales distintos: la cuenta binaria y el avance de una FSM. La secuencia observada en placa fue consistente con las simulaciones, mientras que el reloj de 100 MHz fue mantenido como único dominio en ambos diseños. La división temporal fue lograda mediante clock enable, sin utilizar una salida lógica como reloj adicional.

La separación entre control y datapath fue validada al mantenerse el código de operación independiente de los operandos. Para A=5 y B=3 fueron obtenidos los resultados 1, 7, 6 y 8; en la suma 15+3 fue preservado el quinto bit, con lo que fue evitada la pérdida del acarreo. El reset de la FSM fue distinguido del borrado de datos: se regresó a AND, pero los operandos permanecieron definidos por los switches.

La cobertura de prueba fue mejorada al detectarse que un reset aplicado sobre una cuenta ya nula no demostraba el comportamiento buscado. También fue comprobado que una verificación exhaustiva podía ser correcta y, a la vez, poco clara visualmente; por ello fue añadida una prueba estable de acarreo. La reorganización de los DIP redujo la ambigüedad de lectura de entradas. Como mejora posterior fue identificada la incorporación de sincronización de switches y un registro sistemático de las pruebas físicas. Los márgenes temporales internos positivos y el DRC sin violaciones respaldaron la implementación, aunque no sustituyeron la comprobación del hardware ni una medición instrumental del período.
